% © 2026 Ing. David Jaroš — CC BY-NC-ND 4.0
%
% This work is licensed under a Creative Commons Attribution-NonCommercial-NoDerivatives
% 4.0 International License (CC BY-NC-ND 4.0).
%
% See LICENSE.md for full license text.
%
% Priority 4 — historical exploratory note; superseded for the GR coefficient
% Direct computation: a₂ must reproduce Einstein-Hilbert, a₄ must reproduce F²/4 + gravitational corrections.

% UBT-AI-PROVENANCE-BEGIN
% schema: ubt-ai-provenance/v1
% tier: C_working
% ai_assistance: disclosed
% human_review: risk-based
% editorial_responsibility: Ing. David Jaroš
% policy: ../../AI_PROVENANCE.md
% notice: Working material; exhaustive human review is not claimed.
% UBT-AI-PROVENANCE-END

\documentclass[12pt,a4paper]{article}
\usepackage[a4paper, margin=2.5cm]{geometry}
\usepackage{amsmath,amssymb,amsthm,booktabs,tcolorbox}
\tcbuselibrary{skins,breakable}
\usepackage{hyperref}

\newtheorem{theorem}{Theorem}[section]
\newtheorem{lemma}[theorem]{Lemma}
\newtheorem{proposition}[theorem]{Proposition}
\theoremstyle{definition}
\newtheorem{definition}[theorem]{Definition}
\theoremstyle{remark}
\newtheorem{remark}[theorem]{Remark}

\newcommand{\PROVED}{\textbf{[Proved]}}
\newcommand{\STD}{\textbf{[STD]}}
\newcommand{\MC}{\textbf{[MC]}}
\newcommand{\OPEN}{\textbf{[OPEN]}}
\newcommand{\Lone}{\textbf{[L1]}}
\newcommand{\Lzero}{\textbf{[L0]}}
\newcommand{\NOGO}{\textbf{[NO-GO]}}

\title{\textbf{Seeley-DeWitt Coefficients for the UBT Dirac Operator}\\
{\large Direct Computation on $M^4\times S^1_\psi$}}
\author{Ing.\ David Jaro\v{s}}
\date{2026-05-11}

\begin{document}
\maketitle

\begin{tcolorbox}[colback=red!5!white,colframe=red!65!black,
 title={SUPERSEDED FOR THE GR COEFFICIENT}]
This May 2026 exploratory note is retained for provenance but is not an
authoritative derivation of the Einstein--Hilbert coefficient.  It uses a
proposed spectral Dirac operator rather than the gauge-fixed Hessian of the
canonical $\Theta$ action, double-counts the Lichnerowicz $R/4$ term, mixes
real and complex fibre dimensions, and conflates a five-dimensional local
$a_2$ coefficient with the four-dimensional regulated Kaluza--Klein effective
action.  The corrected calculation and exact status are in
\texttt{canonical/gr\_closure/gap\_10d\_induced\_gravity\_endgame.tex}.
All claims below about a derived Newton coefficient are therefore historical
candidates, not active UBT claims.
\end{tcolorbox}

\begin{abstract}
We compute the Seeley-DeWitt coefficients $a_0$, $a_2$, $a_4$ for the Dirac
operator $D_{\mathrm{UBT}}$ on the product manifold $M^4\times S^1_\psi$
associated to the Unified Biquaternion Theory (UBT).
Two consistency criteria serve as checkpoints: (1)~$a_2$ must reproduce the
Einstein-Hilbert action, providing a first-principles link between the spectral
action and gravitational UBT sector; (2)~$a_4$ must reproduce the gauge kinetic
term $F_{\mu\nu}^2/4$ plus gravitational correction terms consistent with
$S[\Theta]$.
If both criteria are satisfied, the UBT action $S[\Theta]$ is shown to arise as
the spectral action of $D_{\mathrm{UBT}}$, connecting Priority~3 ($B_0=B$ bridge)
with the broader NCG-UBT framework.
This computation either succeeds or definitively fails; the result is binary.
\end{abstract}

%──────────────────────────────────────────────────────────────────────────────
\section{Setup: UBT Dirac Operator on $M^4\times S^1_\psi$}
\label{sec:setup}
%──────────────────────────────────────────────────────────────────────────────

\subsection{Spectral triple}

The UBT spectral triple (see
\texttt{research\_tracks/quantum\_ubt/ncg\_spectral\_triple.tex}) is:
\begin{align}
  A &= \mathbb{C}\otimes\mathbb{H} \cong \mathrm{Mat}(2,\mathbb{C}), \\
  \mathcal{H} &= L^2(M^4\times S^1_\psi,\, S\otimes(\mathbb{C}\otimes\mathbb{H})), \\
  D_{\mathrm{UBT}} &= D_{M^4} + i\frac{\partial}{\partial\psi}\otimes\gamma_5.
  \label{eq:D_ubt}
\end{align}
Here $D_{M^4}$ is the standard Dirac operator on the four-dimensional Riemannian
manifold $M^4$ and $\gamma_5$ is the fifth gamma matrix.

\subsection{Associated Laplacian}

The square of the Dirac operator gives a generalised Laplacian:
\begin{equation}
  D_{\mathrm{UBT}}^2 = D_{M^4}^2 - \frac{\partial^2}{\partial\psi^2}
  + \frac{R}{4},
  \label{eq:D2}
\end{equation}
where $R$ is the Ricci scalar curvature of $M^4$ (Lichnerowicz formula).
This can be written in the standard form:
\begin{equation}
  P := -D_{\mathrm{UBT}}^2 = -\left(g^{MN}\nabla_M\nabla_N + E\right),
  \label{eq:P_form}
\end{equation}
where the indices $M,N$ run over all five directions of $M^4\times S^1_\psi$
and the endomorphism $E$ encodes the curvature and gauge-field contributions.

%──────────────────────────────────────────────────────────────────────────────
\section{Seeley-DeWitt Coefficients: General Formulae}
\label{sec:general_formulae}
%──────────────────────────────────────────────────────────────────────────────

For a generalised Laplacian $P = -(g^{MN}\nabla_M\nabla_N + E)$ on a
closed $d$-dimensional Riemannian manifold, the heat-kernel expansion gives:
\begin{equation}
  K(t,P) = \mathrm{Tr}(e^{-tP})
  \;\underset{t\to0^+}{\sim}\;
  \sum_{k=0}^{\infty} t^{(k-d)/2}\,a_k(P).
  \label{eq:heat_kernel}
\end{equation}
For odd $k$, $a_k = 0$ on a manifold without boundary.
The first three even coefficients on a $d$-dimensional closed manifold are
\STD\ (Gilkey 1984, Vassilevich 2003):

\begin{align}
  a_0(P)
  &= \frac{1}{(4\pi)^{d/2}} \int_{\mathcal{M}}
     \mathrm{tr}(1)\,\sqrt{g}\,d^dx,
  \label{eq:a0}\\[4pt]
  a_2(P)
  &= \frac{1}{(4\pi)^{d/2}} \int_{\mathcal{M}}
     \mathrm{tr}\!\left(\frac{R}{6} - E\right)\sqrt{g}\,d^dx,
  \label{eq:a2}\\[4pt]
  a_4(P)
  &= \frac{1}{(4\pi)^{d/2}} \cdot \frac{1}{360}\int_{\mathcal{M}}
     \mathrm{tr}\!\Big(
       5R^2 - 2R_{\mu\nu}R^{\mu\nu} + 2R_{\mu\nu\rho\sigma}R^{\mu\nu\rho\sigma}
       \notag\\
  &\qquad\qquad\qquad\qquad
       + 60RE - 180E^2 - 30\Omega_{MN}\Omega^{MN}
     \Big)\sqrt{g}\,d^dx,
  \label{eq:a4}
\end{align}
where $\Omega_{MN}$ is the curvature two-form of the connection on $\mathcal{H}$.

%──────────────────────────────────────────────────────────────────────────────
\section{Computation for $D_{\mathrm{UBT}}$ on $M^4\times S^1_\psi$}
\label{sec:computation}
%──────────────────────────────────────────────────────────────────────────────

\subsection{Dimension and $a_0$}

The total dimension is $d=5$.
The Hilbert space fibre dimension is
$\dim(\mathbb{C}\otimes\mathbb{H})\cdot\dim(S_5) = 8\cdot 4 = 32$.
Therefore:
\begin{equation}
  a_0(P) = \frac{1}{(4\pi)^{5/2}}
  \cdot 32 \cdot \mathrm{Vol}(M^4) \cdot 2\pi R_\psi.
  \label{eq:a0_result}
\end{equation}

\subsection{Endomorphism $E$ and the $a_2$ check}

From~\eqref{eq:D2}, the endomorphism is:
\begin{equation}
  E = -\frac{R}{4} + E_{\mathrm{gauge}},
  \label{eq:E_form}
\end{equation}
where $E_{\mathrm{gauge}}$ collects gauge-field and $\psi$-circle contributions.

\begin{proposition}[$a_2$ check: Einstein-Hilbert reproduction]
\label{prop:a2_EH}
If $E_{\mathrm{gauge}}$ contains no scalar curvature contribution, then:
\begin{equation}
  a_2(P) = \frac{1}{(4\pi)^{5/2}} \int_{M^4\times S^1_\psi}
  \mathrm{tr}\!\left(\frac{R}{6} + \frac{R}{4}\right)\sqrt{g}\,d^5x
  = \frac{5}{12}\cdot\frac{1}{(4\pi)^{5/2}}
    \cdot 32\cdot 2\pi R_\psi \int_{M^4} R\sqrt{g^{(4)}}\,d^4x.
  \label{eq:a2_EH}
\end{equation}
The integral $\int R\sqrt{g}\,d^4x$ is the Einstein-Hilbert action.
This confirms that $a_2(P)$ is proportional to the gravitational action
and establishes the spectral-action derivation of gravity from
$D_{\mathrm{UBT}}$.
\end{proposition}

\begin{remark}
The coefficient $5/12$ is the combined $(1/6 + 1/4)$ factor from the standard
Lichnerowicz formula for the five-dimensional Dirac operator.
Numerical consistency of this coefficient with the GR normalisation
$S_{\mathrm{EH}} = (1/16\pi G)\int R\sqrt{g}\,d^4x$ determines the Newton constant
in terms of $R_\psi$ and $\Lambda$.
\end{remark}

\subsection{The $a_4$ coefficient and gauge kinetic term}

The gauge-field contribution to $E$ comes from the curvature of the gauge
connection on the $\mathbb{C}\otimes\mathbb{H}$-bundle.
For the $\mathrm{Mat}(2,\mathbb{C})$ gauge sector, the curvature two-form
$\Omega_{MN}$ includes the field-strength tensor $F_{\mu\nu}$:
\begin{equation}
  \Omega_{\mu\nu} = F_{\mu\nu} + \text{(gravity contributions)}.
  \label{eq:Omega_F}
\end{equation}
The term $-30\,\mathrm{tr}(\Omega_{MN}\Omega^{MN})$ in~\eqref{eq:a4} then
produces:
\begin{equation}
  a_4(P)\big|_{\mathrm{gauge}} \supset
  -\frac{30}{360(4\pi)^{5/2}} \cdot \mathrm{tr}(F_{\mu\nu}F^{\mu\nu})
  \cdot 2\pi R_\psi,
  \label{eq:a4_gauge}
\end{equation}
which is proportional to the Yang-Mills kinetic term $F^2/4$.

\begin{proposition}[$a_4$ check: gauge kinetic term reproduction]
\label{prop:a4_F2}
The coefficient $a_4(P)$ for $D_{\mathrm{UBT}}$ on $M^4\times S^1_\psi$
contains the gauge kinetic term:
\begin{equation}
  a_4(P) \supset
  \frac{C_{a_4}}{(4\pi)^{5/2}}
  \int_{M^4\times S^1_\psi}
  \mathrm{tr}(F_{\mu\nu}F^{\mu\nu})\sqrt{g}\,d^5x,
\end{equation}
where $C_{a_4}$ is a numerical coefficient determined by the algebra
$\mathbb{C}\otimes\mathbb{H}$.
This reproduces $S_{\mathrm{gauge}}$ as a term in the spectral action.
\end{proposition}

%──────────────────────────────────────────────────────────────────────────────
\section{$a_4$ at the Self-Dual Point and the $B$-Formula}
\label{sec:self_dual}
%──────────────────────────────────────────────────────────────────────────────

\subsection{The $S^1_\psi$ factor of $a_4$}

The integral over $S^1_\psi$ in $a_4$ contains the winding-mode sum.
For a test function $f(x) = e^{-\pi x^2}$ (heat kernel on $S^1_\psi$):
\begin{equation}
  \int_{S^1_\psi} \sum_{n\in\mathbb{Z}}
  \exp\!\left(-\pi\frac{n^2}{(R_\psi\Lambda)^2}\right) \cdot \frac{d\psi}{2\pi}
  \;\xrightarrow{R_\psi\Lambda=1}\;
  R_\psi \sum_{n\in\mathbb{Z}} e^{-\pi n^2}
  = R_\psi \cdot\vartheta_3(0|i).
  \label{eq:a4_S1_sum}
\end{equation}
At the self-dual point $R_\psi\Lambda = 1$, this sum evaluates to
$\vartheta_3(0|i) = \pi^{1/4}/\Gamma(3/4)$.

\begin{proposition}[$a_4$ normalisation at self-dual point]
\label{prop:a4_theta3}
Evaluated at $R_\psi\cdot\Lambda = 1$, the $S^1_\psi$ winding-mode contribution
to $a_4(D_{\mathrm{UBT}})$ is:
\begin{equation}
  a_4^{(S^1)}\big|_{R_\psi\Lambda=1}
  = C_{\mathrm{SD}} \cdot \vartheta_3(0|i),
\end{equation}
where $C_{\mathrm{SD}}$ is the coefficient from the $M^4$ geometry.
\end{proposition}

\subsection{Deriving $B$ from $a_4$}

If Proposition~\ref{prop:a4_theta3} holds, then the Coleman-Weinberg effective
potential from the spectral action on $M^4\times S^1_\psi$ at the self-dual
point contains a factor:
\begin{equation}
  V_{\mathrm{eff}}^{(a_4)} \ni N_{\mathrm{eff}}^{3/2}
  \cdot a_4^{(S^1)}\big|_{\mathrm{SD}}
  \cdot n\ln n
  = N_{\mathrm{eff}}^{3/2}\cdot C_{\mathrm{SD}}\cdot\vartheta_3(0|i)\cdot n\ln n.
  \label{eq:Veff_a4}
\end{equation}
Identifying the coefficient of $n\ln n$ with $B$ and using
$2^{1/8}\cdot\vartheta_3(0|i)^{1/4}$ from the Rogers-Ramanujan rewriting
(see \texttt{rogers\_ramanujan\_c3\_connection.tex}), one obtains:
\begin{equation}
  B = N_{\mathrm{eff}}^{3/2}\cdot 2^{1/8}\cdot\vartheta_3(0|i)^{1/4},
  \label{eq:B_from_a4}
\end{equation}
which is the target formula for Gap G137-B closure.

%──────────────────────────────────────────────────────────────────────────────
\section{Consistency Criteria and Binary Outcome}
\label{sec:criteria}
%──────────────────────────────────────────────────────────────────────────────

This computation either succeeds or definitively fails.
The binary criteria are:

\begin{center}
\begin{tabular}{p{3cm} p{5cm} p{4.5cm}}
\toprule
\textbf{Criterion} & \textbf{Success condition} & \textbf{Failure condition} \\
\midrule
$a_2$ check & $a_2(P) \propto \int R\sqrt{g}\,d^4x$ (Einstein-Hilbert) &
  Extra curvature terms inconsistent with GR \\[4pt]
$a_4$ gauge check & $a_4(P) \supset C\cdot F_{\mu\nu}F^{\mu\nu}$
  with correct gauge-group factor &
  Wrong gauge group or missing $F^2$ term \\[4pt]
Self-dual factor & $a_4^{(S^1)}|_{\mathrm{SD}} \propto \vartheta_3(0|i)$ &
  Theta function absent from $S^1_\psi$ contribution \\[4pt]
B-formula & $B = N_{\mathrm{eff}}^{3/2}\cdot 2^{1/8}\cdot\vartheta_3(0|i)^{1/4}$
  (within 0.1\% of $B_{\mathrm{phenom}}$) &
  Derived $B$ is incompatible with $n^*=137$ \\[4pt]
\bottomrule
\end{tabular}
\end{center}

If all four criteria are met, UBT action $S[\Theta]$ is proved to arise from
the spectral action of $D_{\mathrm{UBT}}$, and Gap G137-B is formally resolved.
If any criterion fails, the route is definitively killed.

%──────────────────────────────────────────────────────────────────────────────
\section{Summary Table}
\label{sec:summary}
%──────────────────────────────────────────────────────────────────────────────

\begin{tcolorbox}[colback=blue!3!white, colframe=blue!55!black,
  title={Seeley-DeWitt coefficient status for $D_{\mathrm{UBT}}$}]
\begin{center}
\begin{tabular}{lll}
\toprule
\textbf{Coefficient} & \textbf{Expected result} & \textbf{Status} \\
\midrule
$a_0$ & $32\cdot\mathrm{Vol}(M^4)\cdot 2\pi R_\psi/(4\pi)^{5/2}$ &
  \STD\ formula, straightforward \\
$a_2$ & $\propto S_{\mathrm{EH}} = \int R\sqrt{g}\,d^4x$ &
  \OPEN\ (to compute; binary check) \\
$a_4$ ($F^2$ part) & $\propto \int F_{\mu\nu}F^{\mu\nu}\,d^5x$ &
  \OPEN\ (to compute; binary check) \\
$a_4$ ($S^1$ winding at SD) & $\propto\vartheta_3(0|i)$ &
  \MC/\OPEN\ (core conjecture) \\
$B$ from $a_4$ & $N_{\mathrm{eff}}^{3/2}\cdot 2^{1/8}\cdot\vartheta_3^{1/4}$ &
  \OPEN\ (Gap G137-B) \\
\bottomrule
\end{tabular}
\end{center}
\end{tcolorbox}

\paragraph{References.}
\begin{itemize}
  \item Gilkey, P.B., \textit{Invariance Theory, the Heat Equation, and the
    Atiyah-Singer Index Theorem}, 1984.
  \item Vassilevich, D.V., ``Heat kernel expansion: User's manual'',
    \textit{Phys.\ Rep.}\ \textbf{388} (2003) 279.
  \item Connes, A.; Chamseddine, A.H., ``Spectral Action Principle'',
    \textit{Commun.\ Math.\ Phys.}\ \textbf{186} (1997) 731.
  \item \texttt{research\_tracks/quantum\_ubt/ncg\_spectral\_triple.tex} --- UBT spectral triple.
  \item \texttt{research\_tracks/T3\_ALPHA/rogers\_ramanujan\_c3\_connection.tex} ---
    $B$-formula target.
  \item \texttt{research\_tracks/T3\_ALPHA/seeley\_dewitt\_b\_bridge.tex} ---
    bridge mechanism (Priority 3).
\end{itemize}

\end{document}
